Primero que todo, quiero agradecer al profesor Jorge Vera, por confiar en mí y en mis habilidades. Le agradezco su apoyo al permitirme elegir el tema, lo cual me llevó a conocer un mundo complejo e interesante y a realizar un aporte a algo tan significativo como la salud pública. Esta experiencia ha renovado mi deseo de aprendizaje y crecimiento personal.

También quiero agradecer al equipo de DGR del Servicio Metropolitano Sur, en especial a Diego Rojas, María José Muñoz y Ricardo Ahumada, por su constante disposición para ayudarme y motivarme en este tema. Y a los doctores Francisco Álvarez, Francisco López y Laura Itriago, por su guía en el ámbito médico.

Quiero agradecer a mis familiares, especialmente a mis padres Javier y Gloria, por darme la capacidad de creer en mi mismo. A mis hermanos Banca, Victoria y Josefina, por hacerme parte de su vida y compartir nuestros éxitos y desgracias juntos. También a mis abuelos, tíos y primos, que jamás dejaron de preguntarme y alentarme en la confección de esta tesis. Un agradecimiento especial a mis abuelos Eladio y Josefina, y a Gloria Cartes, por acogerme en su casa durante el verano donde pude dedicarme a escribir la mayor parte de este trabajo.

Finalmente, quiero agradecer a mis amigos, con quienes tengo mis más preciados recuerdos en estos años de estudio. A los del colegio, de los que jamás me sentí alejado, a pesar de haber tomado caminos distintos. Y a mis compañeros de universidad, los Peppas, quienes hicieron todo lo académico más difícil de lo que tenía que ser, pero todo el resto infinitamente más divertido. Finalmente, a Rai y Opa, gracias por acompañarme a enfrentar juntos el desafío de estudiar ingeniería en esta universidad.